\documentclass[tikz,border=10pt]{standalone}
\usepackage[no-math]{fontspec}
\usepackage[brazil]{babel}
\usepackage{amsmath,amssymb}
\IfFileExists{C:/Windows/Fonts/Lato_Regular.ttf}{%
  \setmainfont{Lato}[
    Path = C:/Windows/Fonts/,
    UprightFont = Lato_Regular.ttf,
    BoldFont = Lato_Bold.ttf,
    ItalicFont = Lato_Italic.ttf,
    BoldItalicFont = Lato_BoldItalic.ttf
  ]
  \setsansfont{Lato}[
    Path = C:/Windows/Fonts/,
    UprightFont = Lato_Regular.ttf,
    BoldFont = Lato_Bold.ttf,
    ItalicFont = Lato_Italic.ttf,
    BoldItalicFont = Lato_BoldItalic.ttf
  ]
}{%
  \RequirePackage[default]{lato}
}
\renewcommand{\familydefault}{\sfdefault}

\usepackage{xcolor}
\usepackage{tikz}
\usepackage{pgfplots}
\pgfplotsset{compat=1.18}
\usetikzlibrary{arrows.meta,positioning,shapes.geometric,fit,backgrounds,calc,decorations.pathreplacing,patterns}

% Paleta institucional UEFS / Bioengenharia de Solos
\definecolor{bsNavy}{HTML}{1F4268}
\definecolor{bsNavyDark}{HTML}{152E4A}
\definecolor{bsGold}{HTML}{B38A3D}
\definecolor{bsGoldDark}{HTML}{8C6927}
\definecolor{bsTeal}{HTML}{0E7C7B}
\definecolor{bsForest}{HTML}{2D6A4F}
\definecolor{bsAlert}{HTML}{9E2A2B}
\definecolor{bsEarth}{HTML}{8C5A3C}
\definecolor{bsWater}{HTML}{2E86AB}
\definecolor{bsInk}{HTML}{26313D}
\definecolor{bsMuted}{HTML}{5B6573}

\definecolor{bgNavy}{HTML}{EDF2F7}
\definecolor{bgGold}{HTML}{F8F3E9}
\definecolor{bgTeal}{HTML}{EBF6F5}
\definecolor{bgForest}{HTML}{EEF5F0}
\definecolor{bgAlert}{HTML}{FDF2F2}

\tikzset{
  >=Stealth,
  bsBox/.style={
    draw=#1, fill=white, rounded corners=5pt, line width=1.15pt,
    inner sep=7pt, align=center, font=\footnotesize\sffamily, text=bsInk
  },
  bsTint/.style={
    draw=#1, fill=#1!8!white, rounded corners=5pt, line width=1.15pt,
    inner sep=7pt, align=center, font=\footnotesize\sffamily, text=bsInk
  },
  bsBadge/.style={
    fill=#1, text=white, rounded corners=9pt, inner xsep=8pt, inner ysep=3pt,
    font=\scriptsize\bfseries\sffamily
  },
  bsTitle/.style={font=\normalsize\bfseries\sffamily, text=bsNavy, align=center},
  bsSub/.style={font=\footnotesize\sffamily, text=bsMuted, align=center},
  bsArrow/.style={->, >=Stealth, line width=1.35pt},
  bsDArrow/.style={<->, >=Stealth, line width=1.35pt}
}


\begin{document}
\begin{tikzpicture}

% Título
\node[bsTitle] (title) at (7.3, 6.45)
  {Monitoramento de Pegamento: Índice de Velocidade de Germinação (IVG) e Análise Radicular};
\node[bsSub] (sub) at (7.3, 5.95)
  {Avaliação precoce de emergência em campo e quantificação morfométrica do sistema radicular};

% ================= PAINEL ESQUERDO: PROTOCOLO IVG =================
\begin{scope}[shift={(0,0)}]
  \draw[draw=bsNavy, fill=bgNavy!55, rounded corners=6pt, line width=1.2pt]
    (-0.4, -2.05) rectangle (6.9, 5.35);
  \node[bsBadge=bsNavy, rounded corners=4pt] at (3.25, 5.00) {PROTOCOLO IVG (VIGOR PRECOCE)};

  % Fórmula destacada
  \node[fill=white, draw=bsNavy, rounded corners=5pt, line width=1pt,
        text width=6.1cm, inner sep=6pt, align=center, font=\small\sffamily, text=bsInk]
    (eq) at (3.25, 3.65) {%
      $\displaystyle \mathrm{IVG} = \sum_{i=1}^{k} \frac{G_i}{N_i}$\\[4pt]
      {\scriptsize $G_i$: n\textsuperscript{o} de plântulas/brotações no dia $i$ \quad|\quad $N_i$: dias após o plantio}
    };

  % Etapas do protocolo
  \node[bsBox=bsTeal, text width=2.75cm, minimum height=1.40cm] (iv1) at (1.50, 1.38) {%
    \textbf{1. Contagem Diária}\\[3pt]
    {\scriptsize Registro de brotações\\e emergência em\\parcelas e estacas}
  };
  \node[bsBox=bsNavy, text width=2.75cm, minimum height=1.40cm] (iv2) at (5.00, 1.38) {%
    \textbf{2. Vigor Precoce}\\[3pt]
    {\scriptsize Comparação de\\velocidade de pegamento\\(estacas e sementes)}
  };
  \draw[bsArrow, bsNavy] (iv1) -- (iv2);

  % Mini gráfico ilustrativo de IVG com ampla folga vertical
  \begin{scope}[shift={(0.85, -1.65)}]
    \draw[->, bsMuted, line width=0.8pt] (0, 0) -- (5.0, 0) node[right, font=\scriptsize\sffamily, text=bsMuted] {Dias};
    \draw[->, bsMuted, line width=0.8pt] (0, 0) -- (0, 1.05) node[above=1pt, font=\scriptsize\sffamily, text=bsMuted] {Emerg. (\%)};
    \draw[bsForest, line width=1.4pt] (0, 0) .. controls (1.1, 0.10) and (1.8, 0.85) .. (4.7, 0.92);
    \draw[bsGoldDark, dashed, line width=1.2pt] (0, 0) .. controls (2.1, 0.05) and (3.1, 0.38) .. (4.7, 0.56);
    \node[font=\scriptsize\bfseries\sffamily, text=bsForest] at (3.45, 1.08) {IVG alto (rápido)};
    \node[font=\scriptsize\sffamily, text=bsGoldDark] at (3.75, 0.24) {IVG baixo};
  \end{scope}
\end{scope}

% ================= PAINEL DIREITO: ANÁLISE RADICULAR =================
\begin{scope}[shift={(7.7,0)}]
  \draw[draw=bsForest, fill=bgForest!55, rounded corners=6pt, line width=1.2pt]
    (-0.4, -2.05) rectangle (6.9, 5.35);
  \node[bsBadge=bsForest, rounded corners=4pt] at (3.25, 5.00) {ANÁLISE RADICULAR POR IMAGEM};

  \node[bsBox=bsEarth, text width=6.0cm, minimum height=1.05cm, align=left, inner sep=7pt] (r1) at (3.25, 3.85) {%
    \textbf{1. Lavagem cuidadosa do monolito}\\[1pt]
    {\scriptsize Peneiramento sob jato suave para preservar raízes finas sem ruptura}
  };

  \node[bsBox=bsTeal, text width=6.0cm, minimum height=1.05cm, align=left, inner sep=7pt] (r2) at (3.25, 2.08) {%
    \textbf{2. Escaneamento / fotografia com escala métrica}\\[1pt]
    {\scriptsize Fundo contrastante de alta resolução e régua de calibração espacial}
  };

  \node[bsBox=bsForest, text width=6.0cm, minimum height=1.05cm, align=left, inner sep=7pt] (r3) at (3.25, 0.30) {%
    \textbf{3. Processamento digital de parâmetros}\\[1pt]
    {\scriptsize Comprimento radicular, densidade, \textit{RAR} e distribuição por diâmetro}
  };

  \draw[bsArrow, bsEarth] (r1.south) -- (r2.north);
  \draw[bsArrow, bsTeal] (r2.south) -- (r3.north);

  \node[fill=white, draw=bsForest, rounded corners=4pt, text width=6.1cm, inner sep=5pt,
        font=\scriptsize\bfseries\sffamily, text=bsForest, align=center] at (3.25, -1.25)
    {Saída de projeto: estimativa da coesão adicional das raízes ($\Delta c_r$)};
\end{scope}

\end{tikzpicture}
\end{document}
